\documentclass[11pt]{amsart}

\usepackage[T1]{fontenc}
\usepackage{lmodern}
\usepackage{microtype}
\usepackage{mathtools,amssymb,amsthm}
\usepackage[hidelinks]{hyperref}

\hypersetup{
 pdftitle={A 28-vertex circulant with integral fractional broadcast number
strictly between its multipacking and broadcast numbers}
}

\newtheorem{theorem}{Theorem}
\newtheorem{lemma}[theorem]{Lemma}
\newtheorem{proposition}[theorem]{Proposition}
\theoremstyle{remark}
\newtheorem*{remark}{Remark}
\newtheorem*{problem}{Problem 4 (Teshima)}

\newcommand{\mult}{\operatorname{mp}}
\newcommand{\gb}{\gamma_b}
\newcommand{\gbf}{\gamma_{b,f}}
\newcommand{\ecc}{\operatorname{e}}
\newcommand{\rad}{\operatorname{rad}}
\newcommand{\diam}{\operatorname{diam}}
\newcommand{\Zn}[1]{\mathbb Z_{#1}}

\title[A circulant with integral $\gbf$ strictly between $\mult$ and $\gb$]
{A 28-vertex circulant with integral fractional broadcast number strictly
between its multipacking and broadcast numbers}

\date{}

\subjclass[2020]{05C69, 05C12, 90C05}
\keywords{broadcast domination, fractional broadcast number, multipacking,
circulant graph, linear programming duality}

\begin{document}

\begin{abstract}
Problem~4 of the open-problems section of Teshima's multipacking survey asks
whether some graph $G$ has $\gbf(G)$ an integer with
$\mult(G)<\gbf(G)<\gb(G)$. We exhibit such a graph. For the circulant
$G=C_{28}(1,4)$ on $\Zn{28}$,
\[
 \mult(G)=3<\gbf(G)=4<\gb(G)=5 ,
\]
and the middle term is the integer $4$. Each of the three values is certified by hand: an
explicit primal and an explicit dual of the broadcast linear program pin $\gbf=4$ by
weak duality, a ball-counting bound pins $\gb=5$, and the multipacking number is pinned
between an exhibited $3$-element multipacking and a mod-$7$ tiling obstruction. The dual
is the instantiation at $r=3$ of the vertex-transitive formula that is Theorem~10 of the
survey itself; the only argument here that is not such an instantiation is the mod-$7$
obstruction pinning $\mult=3$. We make no claim that this is the first such witness: we
did not determine the literature status of the problem, and did not read several of the
most directly relevant items.
\end{abstract}

\maketitle

\section{The problem}

Throughout, $G=(V,E)$ is a connected graph, $\ecc(v)$ is the eccentricity of $v$, and
$N_k[v]=\{u: d(u,v)\le k\}$. A \emph{broadcast} on $G$ is a map
$f:V\to\{0,1,\dots,\diam(G)\}$ with $f(v)\le \ecc(v)$ for all $v$; it is
\emph{dominating} if every $u\in V$ lies in $N_{f(v)}[v]$ for some $v$ with $f(v)\ge 1$;
its \emph{cost} is $\sigma(f)=\sum_v f(v)$, and $\gb(G)$ is the least cost of a
dominating broadcast. A set $M\subseteq V$ is a \emph{multipacking} if
\begin{equation}\label{eq:mp}
 |M\cap N_s[v]|\le s
 \qquad\text{for every }v\in V\text{ and every }1\le s\le\diam(G),
\end{equation}
and $\mult(G)$ is the largest size of a multipacking. Finally $\gbf(G)$ is the optimum of
the broadcast linear program
\begin{equation}\label{eq:lp}
 \gbf(G)=\min\Bigl\{\,c\cdot x \;:\; Ax\ge \mathbf 1,\ x\ge 0\,\Bigr\},
\end{equation}
whose columns are the pairs $(v,k)$ with $1\le k\le \ecc(v)$, whose rows are the vertices,
with $A_{u,(v,k)}=1$ iff $d(u,v)\le k$, and with cost $c_{(v,k)}=k$. Linear programming
duality gives $\mult\le \gbf\le\gb$. All of this is verbatim from the source below, whose
notation we keep.

\begin{problem}
Does there exist a graph $G$ with integral $\gbf(G)$ such that
$\mult(G)<\gbf(G)<\gb(G)$?
\end{problem}

\noindent The statement is quoted verbatim, in the survey's own notation, from the fourth of
the seven problems of the \emph{Open Problems} subsection of the \emph{Conclusions} of
\cite{Teshima2014}; it carries no citation there.

\section{The witness}

Let $G=C_{28}(1,4)$ be the circulant on $\Zn{28}$ in which $i$ is joined to $i\pm1$ and to
$i\pm4$. It is a connected, $4$-regular, vertex-transitive graph with $28$ vertices and
$56$ edges; in \texttt{graph6} it is the $64$-byte line
\begin{center}\small
\verb|[hdHGcH@GC_H?H?C_@G?H??c?@G?@G??c??H??@G??C_??H_??HO??Cc??@K_??H|
\end{center}
Breadth-first search from $0$ gives the levels
\[
 \{0\},\quad
 \{1,4,24,27\},\quad
 \{2,3,5,8,20,23,25,26\},\quad
 \{6,7,9,12,16,19,21,22\},
\]
\[
 \{10,11,13,15,17,18\},\quad \{14\},
\]
so $\ecc(v)=5$ for every $v$, $\rad(G)=\diam(G)=5$, and, writing $n_k=|N_k[v]|$,
\begin{equation}\label{eq:profile}
 (n_1,n_2,n_3,n_4,n_5)=(5,13,21,27,28).
\end{equation}
Because $\rad=\diam$ here, the convention in \eqref{eq:mp} and the variant that ranges
$s$ only up to $\rad(G)$, common elsewhere in the literature, agree on $G$.

\begin{theorem}\label{thm:main}
For $G=C_{28}(1,4)$,
\[
 \mult(G)=3<\gbf(G)=4<\gb(G)=5,
\]
and $\gbf(G)=4$ is an integer. So $G$ has the property Problem~4 asks for.
\end{theorem}

The proof is four short certificates. Only \eqref{eq:profile}, the level lists above and
the single set
\begin{equation}\label{eq:T}
 T:=\{v: d(0,v)\ge 4\}=\{10,11,13,14,15,17,18\}=14+\{0,\pm1,\pm3,\pm4\}
\end{equation}
are used, with integer arithmetic throughout except for the two fixed rationals $1/21$
and $1/7$ in (A) and (B).

\subsection*{(A) $\gbf(G)\le4$}
Put $x_{(v,3)}=\tfrac1{21}$ for all $28$ vertices $v$ and $x=0$ on every other column. The
column $(v,3)$ is legal because $3\le \ecc(v)=5$. Each $u\in V$ satisfies
$(Ax)_u=|N_3[u]|/21=n_3/21=1$, so $Ax=\mathbf 1\ge\mathbf 1$, and the cost is
$28\cdot 3/21=4$.

\subsection*{(B) $\gbf(G)\ge4$}
Put $y_v=\tfrac17$ for all $v$. For each column $(v,k)$, $1\le k\le5$,
\[
 \sum_{u\in N_k[v]}y_u=\frac{n_k}{7}
 =\frac57,\ \frac{13}7,\ 3,\ \frac{27}7,\ 4
 \quad\text{against}\quad
 c_{(v,k)}=k=1,2,3,4,5 ,
\]
so $y$ is dual feasible, with equality exactly at $k=3$; its objective is
$\sum_v y_v=28/7=4$. Weak duality alone now gives $\gbf(G)\ge4$, and with (A),
$\gbf(G)=4$.

\smallskip
\begin{sloppypar}
This dual is not ours. It is the instantiation at $r=3$ of the vertex-transitive formula
$\gbf(G)=n\cdot\min_{1\le r\le\rad(G)} r/n_r$, which is Theorem~10 of the source survey
itself \cite{Teshima2014}, where it is credited to R.~C.~Brewster and L.~Duchesne \cite{BD}.
Here $\min(1/5,2/13,1/7,4/27,5/28)=1/7$, attained at $r=3$, and $28\cdot\tfrac17=4$.
\end{sloppypar}

\subsection*{(C) $\gb(G)=5$}
Upper bound: $f(0)=5$ and $f\equiv0$ elsewhere is a legal broadcast, since $5\le\ecc(0)=5$
and $5\le\diam(G)$, and it dominates because $N_5[0]=V$. So $\gb(G)\le5$.

Lower bound, by counting only. A dominating broadcast of cost $c$ assigns positive values
$k_1,\dots,k_m$ with $\sum k_i=c$ and $1\le k_i\le5$, and covers at most
$\sum_i n_{k_i}$ vertices. At $c=4$ the multisets are
\[
 [4]\mapsto 27,\quad
 [3,1]\mapsto 26,\quad
 [2,2]\mapsto 26,\quad
 [2,1,1]\mapsto 23,\quad
 [1,1,1,1]\mapsto 20,
\]
with maximum $27<28$; and the same recursion gives maxima $5,13,21$ at $c=1,2,3$, all
below $28$. So no broadcast of cost at most $4$ dominates, and $\gb(G)=5$.

\subsection*{(D) $\mult(G)=3$}
\emph{Lower bound.} $M=\{0,6,15\}$ has pairwise distances $3,3,4$. Over all $28$ vertices
$v$, $\max_v|M\cap N_1[v]|=1$ and $\max_v|M\cap N_2[v]|=2$; concretely
$N_2[0]\cap N_2[6]=\{1,2,3,4,5,8,26\}$, which is disjoint from
$N_2[15]=\{7,10,11,\dots,20,23\}$, so no vertex has all three of $M$ within distance $2$.
For $s\ge3$ the condition \eqref{eq:mp} is vacuous since $|M|=3$. Hence $M$ is a
multipacking and $\mult(G)\ge3$.

\emph{Upper bound.} Suppose $M$ is a multipacking with $|M|=4$. Summing the $s=3$ case of
\eqref{eq:mp} over all $v$ and exchanging the order of summation,
\[
 \sum_{v\in V}|M\cap N_3[v]|=\sum_{u\in M}|N_3[u]|=4\cdot 21=84
 \le 28\cdot 3=84 .
\]
Equality is forced, so $|M\cap N_3[v]|=3$ for \emph{every} $v$: each vertex $v$ has exactly
one $u\in M$ with $d(u,v)\ge4$. Writing $F_u=\{v:d(u,v)\ge4\}$, the four sets $F_u$
($u\in M$) therefore partition $\Zn{28}$, and $4\cdot 7=28$ is consistent. By
vertex-transitivity of the circulant, $F_u=u+T$ with $T$ as in \eqref{eq:T}.

Reduce modulo $7$. Let $t_j=|\{\tau\in T:\tau\equiv j\bmod 7\}|$, so
\[
 (t_0,\dots,t_6)=(1,1,0,2,2,0,1),
\]
and let $m_j=|\{u\in M: u\equiv j\bmod 7\}|$, so $\sum_j m_j=4$. Each residue class of
$\Zn{28}$ modulo $7$ contains exactly $4$ elements, and the tiling says every $x\in\Zn{28}$
has exactly one representation $x=u+\tau$ with $u\in M$, $\tau\in T$; counting such pairs
inside a fixed class gives
\begin{equation}\label{eq:conv}
 (m*t)(j):=\sum_{a\in\Zn7} m_a\,t_{j-a}=4
 \qquad\text{for every } j\in\Zn 7 .
\end{equation}
Now $t(x)=1+x+2x^3+2x^4+x^6$ is monic of degree $6$, and so is
$\Phi_7(x)=1+x+\dots+x^6$, which is irreducible over $\mathbb Q$. If $t$ vanished at a
primitive $7$th root of unity $\zeta$ then $\Phi_7\mid t$, and equality of degrees with
both monic would force $t=\Phi_7$ --- false, since $t$ has coefficient $0$ at $x^2$. So
$t(\zeta)\ne0$ for every nontrivial character of $\Zn7$. The right-hand side of
\eqref{eq:conv} is the constant vector $4$, whose Fourier transform is supported on the
trivial character alone; hence $\widehat m$ vanishes at every nontrivial character, $m$ is
constant, and $7m_0=4$ --- impossible over the integers. Therefore no $4$-element
multipacking exists and $\mult(G)\le3$.

\begin{proof}[Proof of Theorem~\ref{thm:main}]
Combine (A)--(D): $\mult(G)=3$, $\gbf(G)=4$, $\gb(G)=5$. The chain is strict and the middle
term is an integer, which is the property Problem~4 asks for. $G$ is connected,
as the problem requires.
\end{proof}

\section{Scope}

\noindent\textbf{Status and priority.} We make no priority claim, and we do not establish
that Problem~4 was still unanswered: several of the most directly relevant items --- among
them the manuscript \cite{BD} to which the survey credits the formula that certificate (B)
instantiates --- we did not read. What is offered here is an explicit witness with hand
certificates; whether it is the first, an independent one, or a rediscovery we do not know.

\medskip
\noindent\textbf{Minimality.} We do not claim that $G$ has least order among graphs
satisfying Problem~4, and nothing here bounds that order from below.

\medskip
\noindent\textbf{Verification.} The program \texttt{verify.py} accompanying this note
re-derives, in exact integer and rational arithmetic and from the \texttt{graph6} string
and the definitions printed above, every quantity claimed in Theorem~\ref{thm:main}: the
ball profile, the primal and dual feasibility of (A) and (B) at all $140$ constraints, the
counting bound of (C), and both halves of (D) --- the latter three independent ways, by the
polynomial argument above, by an exhaustive scan of all $210$ residue distributions, and by an exhaustive scan of all $2{,}925$ four-element subsets containing $0$ (legitimate because
$\Zn{28}$ acts transitively). It covers this graph and nothing else, and states in its own
output what it does not cover.

\begin{thebibliography}{9}

\bibitem{BD}
R.~C. Brewster and L.~Duchesne,
\emph{Broadcast domination and fractional multipackings},
Manuscript, 2013; cited in \cite{Teshima2014}, and not obtainable.

\bibitem{Teshima2014}
L.~E. Teshima,
\emph{Multipackings in graphs},
\href{https://arxiv.org/abs/1409.8057}{arXiv:1409.8057v1} (2014).

\bibitem{Yang2015}
F.~Yang,
\emph{New results on broadcast domination and multipacking},
M.Sc.\ thesis, University of Victoria, 2015, handle \texttt{1828/6627}.

\bibitem{Yang2019}
F.~Yang,
\emph{Limited broadcast domination},
Ph.D.\ thesis, University of Victoria, 2019, handle \texttt{1828/11444}.

\end{thebibliography}

\end{document}
